\subsection{Background on Interfaces for SSL Models}
In~\cite{shih24_interface}, the authors define a framework for utilizing SFM models by specifying 3 modules: the Upstream~($\mathrm{U}$) model, Interface~($\mathrm{I}$), and Downstream~($\mathrm{D}$) model.
The three components implement the following mappings:
\begin{align}
    \mathrm{U}\left(\cdot \right) & : \mathbb{R}^{T'} \mapsto \mathbb{R}^{L\times T\times D} \\
    \mathrm{I}\left(\cdot \right) & : \mathbb{R}^{L\times T\times D} \mapsto \mathbb{R}^{ T\times D} % \\
    %\mathrm{D}\left(\cdot \right) & : \mathbb{R}^{T\times D} \mapsto % \mathbb{R}^{ T\times D}
\end{align} where $T'$ is the temporal length of an input speech utterance, $T$ is the length in frames after downsampling by the upstream model, $L$ is the number of layers and $D$ is the hidden dimension of the upstream model.
The downstream model takes the output of $\mathrm{I}\left(\cdot \right)$ as input for specific downstream tasks.
% In this framework, the upstream model is initialized from a pretrained speech encoder and remain frozen throughout the downstream task fintuning, while the interface and downstream is trainable.
For the default setting of this framework, the upstream model is initialized from a pretrained speech encoder and remains frozen throughout the downstream task finetuning, while the interface and downstream are trainable\footnote{In their work, they also try finetuning all components. However, we do not include this setting in our work due to computational constraints.}.
% \david{This is not always the case, correct? For example in the interface paper we also evaluate under the full-fine tuning case}
The function of the interface is to aggregate information across all layers of the upstream model, and a good interface should generalize to many downstream tasks while maintaining a modest number of trainable parameters.
\cite{shih24_interface} empirically demonstrated that the widely-used weighted sum interface is susceptible to information collision across layers, making it suboptimal.
To this end,~\cite{shih24_interface} proposed an alternative interface design (HConv), which uses a series of 1D convolution layers applied across the upstream model's layer dimension. HConv demonstrated significant improvements over the weighted sum interface across various speech processing tasks with different upstream models.

\subsection{Extending Interfaces to Model and Layer Fusion}
In this work, we extend the definition of the Interface framework to encapsulate not only fusion across layers, but also fusion across multiple upstream models.
Formally, suppose we have $N$ Speech Foundation Models $\left[\mathrm{U_1},\mathrm{U_2}\dots \mathrm{U_N}\right]$ as our upstream models.
After feeding the input utterance $\mathbf{x}\in \mathbb{R}^{T'}$ into each upstream model separately, we obtain a list of 3-dimensional tensors by extracting all hidden representations of each model. i.e., $\mathbf{h}=\left\{\mathbf{h}_1,\mathbf{h}_2,\dots,\mathbf{h}_N\right\}$.
Each hiddenstate $\mathbf{h}_i\in \mathbb{R}^{ L_i \times T_i \times D_i}$ is allowed to have different numbers of layers, temporal length, and hidden dimension. 
Under this relaxation of the framework, an interface module jointly performs model and layer fusion by implementing the mapping $\mathrm{I}\left(\cdot \right): \left( 
\mathbb{R}^{L_1\times T_1\times D_1},
\mathbb{R}^{L_2\times T_2\times D_2},
\dots,
\mathbb{R}^{L_N\times T_N\times D_N}
\right) \mapsto \mathbb{R}^{ T\times D}$.



% Notice that here we assume the upstream models have an identical number of layers and hidden dimensions to make our evaluation and comparison easier.
% In fact, this 
% We leave the fusion of models with different sizes to future work.
% Hence, an interface module jointly performs model and layer fusion implements the mapping $\mathrm{I}\left(\cdot \right): \mathbb{R}^{N \times L\times T\times D} \mapsto \mathbb{R}^{ T\times D}$.


%\subsection{Proposed Revision}
To allow fusion for multiple upstream models, we revise the Hierarchical Convolution Interface.
First, we must ensure that all $L_i$ and $T_i$ are identical across models. i.e. $\forall i,j, 1\leq i,j \leq N, L_i=L_j \wedge T_i=T_j$.
While most of the commonly used Speech Foundation Models utilize a framerate of 50Hz and typically use the same number of layers (12 Transformer layers for base models and 24 for Large models), in the case that the number of layers or the framerates differ for some of the upstream models, we upsample the representations of the models to match that of the model with the highest framerate (and similarly with the number of layers), and then apply linear interpolation to both axes.
We denote the matched layer and temporal dimension as $L$ and $T$ respectively.
After matching $L,T$, we have two options for merging the hidden states of each model.
If the hidden dimensions of each model are the same (or interpolation is applied to make it the same), we can directly add all models' hidden state together to obtain a tensor as $\mathbb{R}^{L \times T\times D}$.
The other option is to directly concatenate all models' hidden states to obtain a tensor as $\mathbb{R}^{L \times T\times \left(N \cdot D\right)}$.
We can then insert the merged tensor into the Hierarchical Convolution Interface.
We call the former method ``HConv.'' and latter ``CHConv.''
Note that the merging operations mentioned above possess the characteristic of local similarity on the last dimension of the tensor, which still complies with the assumption of using the Hierarchical Convolution Interface.




% Our proposed method first apply linear interpolation on layer and temporal dimension to unify each encoder's hiddenstates dimensions.
% Practically, the commonly available Speech Foundation Models
% To allow fusion for multiple upstream models,
% We apply either layer-wise addition or concatenation (with a projection layer) to the upstream model layer outputs before applying the Hierarchical Convolution Interface.
% These operations are chosen because they preserve local similarity along the layer dimension, aligning with the assumptions of the Hierarchical Convolution Interface. 
% Notably, concatenation + projection subsumes addition but is more general.

